\section*{Resolución en la conversión analógico digital}

La resolución en la conversión analógico-digital (ADC) describe la capacidad del sistema para discriminar pequeños cambios en una señal analógica continua al convertirla en valores digitales discretos según \cite{NIResolutionDAQ}. Esta resolución está directamente determinada por el número de bits del convertidor, ya que un ADC de $N$ bits divide el rango completo de voltaje de entrada en $2^N$ niveles cuantizados \cite{LibreTextsADC}. Por ejemplo, un ADC de 16 bits proporciona 65\,536 niveles posibles. Desde el punto de vista práctico, la resolución define el menor cambio de voltaje detectable, conocido como LSB (\textit{Least Significant Bit}), que se calcula dividiendo el rango de entrada configurado (por ejemplo $\pm 10$\,V) entre el número total de niveles. En aplicaciones biomédicas, donde las señales suelen ser de baja amplitud, esta característica es crítica para no perder información fisiológica relevante.

En las tarjetas de adquisición de datos (DAQ), como las desarrolladas por National Instruments (NI), la resolución del ADC debe analizarse junto con otros parámetros como el ruido del sistema, la exactitud (\textit{accuracy}), la linealidad integral y diferencial, y la estabilidad del voltaje de referencia según \cite{ADLINKDAQ}. Aunque un ADC tenga una alta resolución teórica, el número efectivo de bits (ENOB) suele ser menor debido a ruido electrónico, interferencias electromagnéticas y errores internos del circuito. En entornos hospitalarios, esto cobra especial importancia porque las mediciones deben cumplir criterios de confiabilidad y repetibilidad exigidos por normativas técnicas y de seguridad.

\section*{Mediciones diferenciales con tarjeta de adquisición de datos}

La medición diferencial es una técnica ampliamente utilizada en DAQ de NI para mejorar la calidad de la señal, especialmente cuando se trabaja con señales biomédicas o sensores ubicados en ambientes con alto nivel de ruido. En este modo, el ADC mide la diferencia de voltaje entre dos entradas $(AI+ y AI-)$, en lugar de medir una entrada respecto a tierra. Al hacer esto, se atenúan de manera significativa los ruidos comunes presentes en ambas líneas, como interferencias de la red eléctrica (50/60 Hz), lo que incrementa la relación señal-ruido y permite aprovechar mejor la resolución del ADC.

La relación entre resolución y medición diferencial es directa: al reducir el ruido y las tensiones parásitas, la resolución efectiva del sistema mejora, ya que los niveles digitales del ADC se utilizan para representar información real de la señal y no variaciones espurias. En tarjetas NI, además, la configuración del rango de entrada es fundamental. Seleccionar un rango lo más cercano posible a la amplitud real de la señal permite maximizar el uso de los niveles del ADC, reduciendo el valor del LSB y mejorando la sensibilidad de la medición, lo cual es esencial en aplicaciones como adquisición de ECG, EMG o sensores de presión invasiva.